%% Compile with XeLaTeX (Overleaf: Menu -> Compiler -> XeLaTeX).
\documentclass[conference]{IEEEtran}
\usepackage{iftex}
\RequireXeTeX   % Overleaf: Menu -> Compiler -> XeLaTeX (pdfLaTeX cannot typeset the Hindi quotations)
\usepackage{fontspec}
\setmainfont{texgyretermes}[Extension=.otf, UprightFont=*-regular, BoldFont=*-bold, ItalicFont=*-italic, BoldItalicFont=*-bolditalic]
% Devanagari font for the Hindi: the bundled file if present (zip upload), else a font Overleaf already has.
\IfFileExists{fonts/NotoSerifDevanagari-Regular.ttf}{%
  \newfontfamily\devfont[Path=fonts/, Script=Devanagari, Scale=0.92]{NotoSerifDevanagari-Regular.ttf}}{%
  \IfFontExistsTF{Noto Serif Devanagari}{%
    \newfontfamily\devfont[Script=Devanagari, Scale=0.92]{Noto Serif Devanagari}}{%
    \newfontfamily\devfont[Script=Devanagari]{FreeSerif}}}
\newcommand{\dev}[1]{{\devfont #1}}
\usepackage{amsmath,amssymb}
\usepackage{booktabs,tabularx,array}
\usepackage[hyphens]{url}
\usepackage[hidelinks]{hyperref}
\newcolumntype{Y}{>{\raggedright\arraybackslash}X}
\renewcommand{\arraystretch}{1.12}
\setlength{\tabcolsep}{3.5pt}
\begin{document}

\title{Assessing Legal Meaning Preservation in English and Hindi Simplifications of Indian Statutes}
\author{\IEEEauthorblockN{Karan Singh (S24CSEU1939), Atharv Arora (S24CSEU2439), Jatin Khanchi (S24CSEU1932)}\IEEEauthorblockA{School of Computer Science Engineering and Technology\\ Bennett University, Greater Noida, India\\ \textit{CSET 346 Natural Language Processing (project PID-08), instructor Dr. Surabhi Solanki}}}
\maketitle
\begin{abstract}
Plain-language versions of statutes widen access to the law, but a simplification that drops a condition, turns an obligation into a permission or alters a time limit misleads the reader while still reading fluently. JUDGEBERT introduced a learned metric for legal meaning preservation (LMP) on French insurance text. We extend it to Indian statutory law in English and Hindi. From the Right to Information Act, 2005 and the Consumer Protection Act, 2019 we build 338 original-simplification pairs in three styles, one carrying a rule-injected legal error of a known type. Pairs are scored on JUDGEBERT's 1-10 rubric by two disclosed large-language-model annotators rather than humans, a deviation we report throughout. Trained on this corpus alone a cross-encoder is unstable (Pearson 0.228 \ensuremath{\pm} 0.288) and not reliably better than the best generic baseline (calibrated LaBSE, 0.101). Pre-training on FrJUDGE first gives 0.576 \ensuremath{\pm} 0.102 for our principal XLM-R system; a post-hoc MuRIL variant matches it (0.577 \ensuremath{\pm} 0.068 over five seeds) and generalises better to Hindi and across acts, while a re-run of the FrJUDGE pre-training shows the transfer depends on that run. Zero-shot English-Hindi transfer stays near zero. A variant-mean lookup using the answer key reaches 0.831, so we also report within-variant correlations. For the optional legal-error type, a MuRIL classifier trained with injected-error augmentation reaches macro-F1 0.789 against 0.509 for rules. Injected negations are handled far better than by embedding baselines, obligation-to-permission flips poorly, and no system passes both sanity checks.
\end{abstract}
\begin{IEEEkeywords}
legal NLP, text simplification, meaning preservation, evaluation metrics, Hindi, multilingual transformers, cross-encoder regression
\end{IEEEkeywords}

\section{Introduction}

Indian statutes are drafted in long, heavily qualified sentences, and the official Hindi versions are no easier to read than the English. Large language models can rewrite such text into plain language, and doing so at scale is attractive for access to justice. The risk is that the rewriting is judged by the wrong yardstick. Generic similarity metrics reward lexical and semantic overlap, so a simplification that preserves almost every word but deletes a proviso scores well, even though the legal consequence has changed.

JUDGEBERT [1] made this concrete. On FrJUDGE, a corpus of 297 French insurance sentence blocks scored by five law students, BERTScore correlates with human LMP judgements at Pearson 0.46 and SBERT at 0.13, while a CamemBERT-v2 cross-encoder trained directly on the rubric reaches 0.74. BERTScore and MeaningBERT over-score the human label on 82\% of test pairs; the learned metric does so on none. The lesson is that legal meaning preservation is a distinct target that has to be learned, not approximated by similarity.

We ask whether that result carries to a different legal system, a different pair of languages and a different kind of statutory text, and we attack the specific weaknesses of the original study. Each gap below maps to a contribution of this paper.

\begin{itemize}
\item JUDGEBERT covers French and a single domain. We cover two Indian central acts from different domains, and English and Hindi within one of them, the Right to Information Act, because the official Hindi text of the other could not be obtained (Section III-A). Where both languages exist, the Hindi statute is an official translation of the very same provision, so the language comparison is controlled rather than confounded by content.
\item FrJUDGE contains one simplification style, so its label distribution is narrow. We add aggressive simplifications, which tend to drop conditions, and perturbed simplifications carrying an additional injected legal error of a known type. The latter is a labelled diagnostic set: it tells us which legal error types a metric is blind to.
\item Krippendorff's alpha on FrJUDGE LMP was 0.10, with annotators splitting into a lenient and a strict group. We therefore treat agreement as a first-class result rather than a footnote, and we report it for whichever annotator population we actually used.
\item FrJUDGE splits are pair-level, so a test original can reappear in training as an identical augmentation pair, and the headline 0.97 Pearson of the augmented model is measured on a test set that contains trivial identical and unrelated pairs. Our splits are grouped by source clause and verified disjoint at clause and text level, and sanity checks are reported separately on holdout provisions used nowhere else.
\item JUDGEBERT reports no out-of-domain test. We train on one act and test on the other.
\end{itemize}

Our contributions are therefore: (i) a bilingual legal LMP corpus with a documented, blinded scoring protocol and an explicit statement of who, or what, produced the labels; (ii) a JUDGEBERT-style multilingual cross-encoder for LMP; (iii) a comparison against generic metrics and, in place of the LLM judge our labelling decision ruled out, a transfer experiment that pre-trains on FrJUDGE before fine-tuning on our splits; (iv) cross-lingual, cross-domain and sanity experiments; (v) a legal-error type classifier evaluated with precision, recall and F1; and (vi) a qualitative error analysis. We also report, in Section VII, two failures of our own pipeline that would have produced confidently wrong numbers, because both are invisible unless specifically tested for.

\section{Related Work}

\subsection{Evaluating simplification}

Simplification quality is usually decomposed into simplicity, fluency and meaning preservation. Meaning preservation is commonly measured by n-gram overlap, by contextual embedding similarity such as BERTScore [3], by sentence-embedding cosine similarity, or by learned or question-based metrics such as QuestEval [9], LENS [10] and MeaningBERT [11]. BLESS [2] benchmarks large language models as sentence simplifiers with zero-shot prompts; we reuse its prompt so that our normal variant is comparable to established practice.

\subsection{Legal meaning preservation}

JUDGEBERT [1] is the direct predecessor of this work. Its rubric, which we adopt unchanged, first places a simplification in a bracket (Accurate 7-10, Imprecise 2-6, Off-track 1) and then subtracts one point per identified error, with a floor of 1; the error types are omission, hallucination, consistency issue and confusion. Its model is a cross-encoder regression over the concatenated sentence pair. Its augmentation adds identical pairs labelled 10 and unrelated pairs labelled 0, the latter mined from Quebec legal texts with ROUGE below 0.25 and BLEU below 25.

\subsection{Indian legal NLP}

Work on Indian legal text has concentrated on judgments: domain-adapted encoders such as InLegalBERT [12], and benchmarks such as IL-TUR [13], target tasks like judgment prediction, rhetorical-role labelling and summarisation. Statutory simplification, and the question of whether a simplification is legally faithful, is comparatively unexplored for Indian law, and we are not aware of a Hindi resource for it.

\subsection{Multilingual encoders}

For the backbone we consider XLM-R [5], which is multilingual but not Indic-focused, and MuRIL [6] and IndicBERT v2 [7], which are trained with Indian languages and transliteration in mind. LaBSE [4] provides a strong language-agnostic sentence embedding and serves as one of our baselines.

\section{Dataset}

\subsection{Sources and provenance}

We use two central acts from different legal domains, each chosen because an official Hindi version exists: the Right to Information Act, 2005 (public information) and the Consumer Protection Act, 2019 (consumer law). India Code is the canonical repository for both, but it was unreachable while this version of the corpus was built, so we used other official government sources and recorded exactly which. Table I lists them.

\begin{table}[!t]
\centering
\caption{Sources, editions and extraction yield}
\label{tab:sources}
\footnotesize
\begin{tabularx}{\columnwidth}{YY>{\centering\arraybackslash}p{33pt}>{\centering\arraybackslash}p{21pt}}
\toprule
\textbf{Document} & \textbf{Edition} & \textbf{Sections} & \textbf{Units} \\
\midrule
RTI Act 2005, English (DoPT) & As on 18 Nov 2025; includes the 2023 amendment & 31 & 109 \\
RTI Act 2005, Hindi (DoPT) & Carries the 2019 amendments; sec. 8(1) is the older wording; text layer unusable, read by OCR (Tesseract 5.4, tessdata\_best hin) & 27 & 101 \\
Consumer Protection Act 2019, English (Gazette) & Act 35 of 2019, 9 Aug 2019; original enactment & 107 & 229 \\
Consumer Protection Act 2019, Hindi & Not used: unavailable when the corpus was built (Bill text as introduced found later) & -- & -- \\
\bottomrule
\end{tabularx}
\end{table}

Two consequences follow and are carried through the rest of the paper. First, the official Hindi text of the Consumer Protection Act could not be obtained, so that act contributes English-only material and the English-Hindi comparison (E2) is confined to the Right to Information Act. Second, the two Right to Information copies are not the same edition: the English consolidation is current to 18 November 2025 and includes the 2023 amendment, whereas the Hindi copy, although it carries the 2019 amendments in its footnotes, retains the older personal-information wording in section 8(1). Provisions affected by this are flagged automatically and dropped at review rather than silently reconciled, because editing one language to match the other would fabricate statutory text.

\subsection{Extraction}

Each PDF is converted to sub-section units, the unit of analysis being a numbered sub-section such as section 7(1), with its marginal or inline heading attached. English and Hindi units of the same provision share a key, which is what makes the bilingual alignment automatic. Three details were needed for real government PDFs and are worth recording, since each cost us a parser bug. Cross-references of the form “sub-section (3) of section 8” must not be read as the start of a new sub-section. Amendment markers, footnote blocks and the arrangement-of-sections table at the front of an act must be stripped, but the word “Repeal” must not be mistaken for the footnote abbreviation “Rep.”. The Gazette printing of the Consumer Protection Act places section headings in the page margin rather than inline, so for that document we extract with a geometry-aware backend that reads the marginal headings and the marginal statute citations separately from the body column, and fails loudly if a numbered section is found without a heading.

\subsection{The Hindi text layer}

The Hindi PDF is not scanned, yet its text layer cannot be used as it stands, and the way it fails is instructive. Extracted with Poppler's pdftotext, 4.9\% of all Devanagari characters emerge as U+FFFD replacement characters: 2461 characters across 100 of 102 units, an irreversible loss. Extracted with PyMuPDF, nothing is lost, but the legacy Mangal glyphs are mapped into the Vedic and Sundanese Unicode blocks (2681 characters drawn from only 48 distinct code points) and the vowel signs appear in visual rather than logical order. In other words the page is readable to a human and systematically wrong to a machine. Our extractor therefore computes a Devanagari ratio and a misplaced-vowel-sign ratio for every Hindi document and refuses the text layer when it fails either test; both routes above are correctly rejected. We report this because a pipeline that silently accepted either extraction would have produced a Hindi corpus of fluent-looking nonsense, and no downstream metric would have revealed it.

The Hindi text used in this work therefore comes from optical character recognition: the pages are rendered at 300 dpi and read with Tesseract 5.4 using the high-quality tessdata\_best Hindi model, which the predecessor literature and our own trials both prefer over the fast model for legal vocabulary. This clears the acceptance test decisively: Devanagari ratio 0.993, misplaced-vowel-sign ratio 0.000 and no replacement characters at all. One systematic OCR defect remains and is worth naming, because it attacks precisely the part of a statute that must not be approximated: the digit 1 is frequently dropped or misread, so section 12 is read as 2, section 13 as 03 and the sub-section marker (1) as an empty pair of brackets. We therefore repair numbering from the sequence rather than from appearance, restoring a number only when the observed token is a recognised one-digit-drop variant of the number the sequence expects, and we do not touch anything else. Every repair is logged.

\subsection{Clause selection}

A provision is eligible if it has one to five sentences, 15 to 90 English words and a Flesch reading-ease score of at most 50, so that the simplification task is non-trivial; if it is not boilerplate (short title, extent, definitions, rule-making powers, repeal and savings); and if its text is not duplicated elsewhere in either act. The duplicate test matters more than it sounds: statutes reuse standard provisos verbatim, and two identical provisions in different splits would leak the test set into training even though the split is grouped by clause. At most three clauses are taken per section, and sampling is weighted towards provisions containing exceptions, overrides, conditions and time limits, which is where legal meaning is most fragile. We oversample by a factor of 1.3 to absorb losses at review.

Every candidate then passes a review step that is the only place where a clause can be accepted. The selection tool pre-computes warning flags (replacement characters, Latin characters inside Hindi text, numbers that differ between the two languages, missing Hindi counterpart, known edition mismatch) and never sets the keep decision itself. In the design this review is performed by a human. Under the project deadline it was performed instead by a rule-based agent script, and we disclose that: rows that are misaligned, fragmentary, contain Latin characters inside Hindi or fall in the section where the two editions differ are dropped, with the reason recorded per row. The only text repairs permitted are the numbering repairs described above, and only where the Hindi number differs from its English counterpart by an OCR-dropped 1; seven such repairs were made and each is logged with the before and after text. A sample of kept rows was then checked against the source pages. No agent at any point edited statutory wording.

The resulting corpus is 78 clauses: 38 from the Right to Information Act, of which 35 have an aligned Hindi counterpart, and 40 from the Consumer Protection Act, which are English-only because its Hindi text was unobtainable. That gives 113 clause-language items and, with three variants each less one perturbation rule that found nothing to edit, 338 pairs. The brief asks for 300 pairs with human judgements; our 338 carry LLM judgements, so that requirement is not met (Section III-D). Table II breaks the corpus down.

\begin{table}[!t]
\centering
\caption{Corpus composition by act, language and variant}
\label{tab:dataset}
\footnotesize
\begin{tabularx}{\columnwidth}{YY>{\centering\arraybackslash}p{29pt}>{\centering\arraybackslash}p{21pt}>{\centering\arraybackslash}p{21pt}>{\centering\arraybackslash}p{21pt}}
\toprule
\textbf{Act} & \textbf{Lang} & \textbf{Clauses} & \textbf{Pairs} & \textbf{Mean words} & \textbf{Mean label} \\
\midrule
CPA & EN & 40 & 119 & 50.2 & 7.40 \\
RTI & EN & 38 & 114 & 53.0 & 7.69 \\
RTI & HI & 35 & 105 & 52.5 & 7.49 \\
All &  & 78 & 338 & 51.9 & 7.53 \\
\bottomrule
\end{tabularx}
\end{table}

\subsection{Simplifications}

For each selected clause and each available language we generate three variants. The normal variant uses the zero-shot BLESS prompt, translated for Hindi. The aggressive variant asks for an explanation to a twelve-year-old in at most two sentences, which reliably sacrifices conditions and therefore populates the low end of the scale that FrJUDGE lacks. The perturbed variant takes a normal simplification and applies one rule-based legal edit to it, balanced across five types: modality (shall or must becomes may, and the reverse), negation, a changed number, deletion of an exception, proviso or override, and narrowed scope. The Hindi rules operate only on clause-final verb groups, so that, for example, karega becomes kar sakega without touching unrelated words that merely contain the same letters. The idiom “as may be prescribed” is explicitly exempted, because the may in it confers no permission. The base simplification is not itself verified to be error-free, so a perturbed pair carries one additional injected error, not necessarily exactly one error relative to the statute: the Hindi normal simplification of RTI section 12(6), for instance, already narrows legislature (vidhan-mandal) to assembly (vidhansabha), and its perturbed copy keeps that slip as well.

The API pipeline we built caps output at 200 tokens for English and 500 for Hindi, rather than JUDGEBERT's 100, but it was not used for this corpus, and one deviation must be stated: JUDGEBERT sampled its simplifications from gpt-4-turbo through an API at temperature 1.0 and top-p 0.9, whereas ours were written by a Claude Opus 5.5 agent following the same prompts, with no sampling parameters to record. The prompt identifier, the model string and a timestamp are stored with every generation, and the temperature and top-p fields are null rather than invented. Our simplifications are consequently likely to be more conservative than temperature-1.0 samples, which plausibly shifts the label distribution upward relative to FrJUDGE. Because no sampling API was involved, no output cap was enforced, and the importer's max\_tokens (null) and finish\_reason (always written as stop) fields are not observations of how generation ended.

\subsection{Scoring protocol and its deviation from human judgement}

This is the point at which the present version of the work departs from our protocol, and we state it plainly rather than in a footnote. The intended design was three human annotators per pair after a shared calibration round. Because of the project deadline, the LMP scores used here were produced by two disclosed large-language-model annotators, both instances of Claude Opus 5.5, each filling its own blinded workbook under the same rubric and conventions written for the human annotators. This does not satisfy the requirement of agreement with human legal judgement, and no claim in this paper should be read as establishing it.

Blinding was implemented mechanically for what each annotator could open, and it is only partial, as the next paragraph makes explicit. Each annotator receives only an opaque pair identifier, in an order shuffled per annotator with items from the same clause separated; the variant name, the injected error type and the answer key are in files the annotators never read, and neither annotator sees the other's workbook. The score is not entered directly: the annotator chooses a bracket and counts errors by type, and the score is recomputed from those entries by our own code, so a rubric change cannot be lost in a spreadsheet formula. Two further constraints follow from using models as annotators. Model families overlap between generation and annotation, so any annotator that also generated the text it scores cannot be called blind, and we disclose that overlap rather than assuming it away. And an LLM judge may not be evaluated against labels from its own model, since that measures self-consistency and not meaning preservation. Because both of our annotators are the same model, that rule removed the LLM-judge baseline entirely: any judge we had access to would have been scored against its own family's labels. We therefore report no LLM-judge result, and record its absence as a consequence of the labelling decision rather than quietly dropping a planned comparison. It is the first experiment to run once human labels exist.

Three further caveats belong here rather than in a footnote, because they bound what the agreement number can mean. First, both annotators are instances of the same model, Claude Opus 5.5, so their agreement is between two instances of one model family and not between two independent judges, and their agreement is likely inflated by shared pretraining. Second, the two annotators are not equally blind. The first was given only the opaque identifiers and the text, and knew the corpus design in the abstract but not which items carried an injected error. The second also wrote the simplifications, so although its workbook hides the variant labels it may recognise its own output and infer which copies were perturbed; it is therefore not blind in practice, and we say so. Third, with two annotators the leave-one-out statistic is simply their Pearson correlation. It is not an upper bound on predicting their mean: each annotator correlates 0.904 and 0.949 with the mean label, which includes its own score. We therefore report it only as inter-LLM agreement and draw no percentage-of-ceiling conclusions from it.

On those terms the agreement is as follows. Krippendorff's alpha [8] on the interval LMP scores is 0.610 and on the nominal brackets 0.503; the Pearson correlation between the two annotators (inter-LLM agreement) is 0.723; and 68 of the 338 pairs fall on the adjudication list, none of which has been adjudicated. These figures are far above the 0.10 that JUDGEBERT reports between its five law students, and the comparison must be read with that asymmetry in mind: two instances of one model agreeing is a much weaker result than five people agreeing, and the gap is evidence about model consistency rather than about the construct being well defined. The annotators are also systematically offset: the mean score of the first is 8.115 and of the second 6.935, so the label, their mean, is in places a numerical compromise between incompatible scoring conventions rather than an adjudicated judgement. The label distribution is bimodal, concentrated at 10 with a tail between 1 and 5; individual scores rarely fall on 6 or 7 because the rubric subtracts from bracket maxima of 10, 6 and 1, although mean labels do (15 pairs at 6 and 7 at 7). The bimodality matters for RMSE. One convention differed between the annotators and we state it rather than smoothing it over: where an inserted or dropped negation reverses an entire operative provision, one annotator scored Off-track, the other Imprecise with a confusion. The label is the mean of the two, as in FrJUDGE, and the adjudication list records every such divergence.

To keep a path to the original claim open, a separate 30-pair validation workbook was prepared for human legal annotators, stratified by language and variant. At the time of writing it has not been scored. No human validation of these labels therefore exists, the correlation between our model labels and human legal judgement is unmeasured, and every result in this paper should be read with that gap in mind. The workbook and the tooling to aggregate it are in the repository, so the measurement is a scoring session away rather than a rebuild.

\subsection{Augmentation, sanity holdouts and splits}

Following JUDGEBERT we add, for each clause and language, an identical pair labelled 10 and an unrelated pair whose partner is a provision with ROUGE-1, ROUGE-2 and ROUGE-L below 0.25 and BLEU below 25. We label unrelated pairs 1 rather than 0, because 1 is the floor of the rubric and an unreachable 0 would distort the regression; a flag reproduces the original choice. These overlap statistics are computed with a Devanagari-aware tokenizer: the widely used rouge\_score package discards non-ASCII text, which would classify every Hindi pair as unrelated and quietly poison the augmentation set.

The sanity holdouts are built only from provisions that are not selected, are not in a section containing a selected clause, and are not used as augmentation partners, so that passing a sanity check cannot be an artefact of memorising training text. Splits are 60/10/30 by clause over three seeds (42 to 44), and every pair derived from one clause, both languages, all three variants and its augmentation pairs, stays within a single split; the splitter asserts this at clause level and at text level. Headline metrics are computed on annotated test pairs only, with augmentation and sanity pairs excluded and reported separately.

\section{Method}

\subsection{LMP model}

The model follows JUDGEBERT: the pair is encoded as original, separator, simplification by a pretrained multilingual encoder, and a linear regression head on the pooled representation predicts the score under mean-squared-error loss, with predictions clipped to the rubric range 1 to 10. We train XLM-R and MuRIL (IndicBERT v2 was planned but not run), each with and without augmentation, over three clause-grouped seeds (42-44). Checkpoints are selected on the annotated development set. The principal system, the FrJUDGE-pretrained XLM-R fine-tuned with augmentation, is fixed in advance as the configuration used for every experiment E1-E4; tables list all systems in a fixed order and never choose the featured system by test score.

Our optimisation settings depart from the published ones and the reason is a property of our corpus, so we state both. JUDGEBERT uses learning rate 5e-5, batch size 16 and early stopping with patience 5; we use learning rate 2e-5, batch size 8, no early stopping, and 30 epochs keeping the best development checkpoint, with linear decay and gradient clipping at norm 1.0 unchanged. The reason is that our development split holds roughly 35 annotated pairs, and on a set that small the development loss is dominated by noise: with patience 5 and with patience 10 the run halted early and left the model predicting the label mean, giving a near-zero correlation. The FrJUDGE reproduction in Section VII shows the same pathology on the original data, where patience 5 collapses one seed in ten, so this is a property of the stopping rule interacting with a small dev set rather than of our corpus alone. We disclose that these settings were chosen after observing that failure, which is a form of selection on the development set and is recorded as a limitation.

\subsection{Baselines}

We compare against four reference points of increasing sophistication. ROUGE-L and BLEU, computed with our Devanagari-aware tokenizer, represent pure overlap. LaBSE [4] cosine similarity represents language-agnostic sentence semantics. Multilingual BERTScore [3] F1 represents contextual token matching and is the strongest generic metric in the original study. A fourth reference point was designed and then deliberately not run: an LLM judge receiving the annotators' rubric verbatim at temperature 0, which would have asked whether a general model with the rubric in its prompt can replace a model trained on it. Section III-F explains why it is unusable here, namely that our labels come from the same model family, so the comparison would have measured self-consistency. The implementation exists and is reported as not run rather than removed. Every remaining baseline is mapped onto the 1-10 scale twice, once naively and once by a linear calibration fitted on the training split only, and both are reported, since an uncalibrated metric can correlate well while having a meaningless RMSE.

\subsection{Decisions and reporting}

For the decision output we map the continuous prediction onto the rubric brackets: at least 7 is Preserved, 2 up to 7 is Partially Preserved, below 2 is Not Preserved. Alongside Pearson and RMSE we report Spearman and MAE, and the over-scoring rate, that is the share of pairs whose prediction exceeds the label, because the practically dangerous failure for this task is a metric that flatters a simplification which has lost a condition. All aggregate numbers are means with standard deviations over the three seeds, and results are broken down by language, by act and by variant.

\section{Experiments}

Five experiments are defined in advance. E1 compares the generic baselines against the trained cross-encoder on the random clause-grouped splits; the LLM judge belonged here and is reported as not run for the reason given in Section III-F. E2 reports per-language results and zero-shot transfer in both directions, training on English and testing on Hindi and the reverse, which is possible precisely because the two languages contain the same provisions; it is confined to the Right to Information Act for the reason given in Section III-A. E3 trains on one act and tests on the other in both directions. E4 runs the sanity checks on the disjoint holdouts, reporting the share of identical pairs scored at least 9.9 and of unrelated pairs scored at most 1.1. E5 is a qualitative analysis of at least fifteen legally difficult cases, grouped by injected error type and by the most important change each annotator recorded, which the perturbed variant makes systematic rather than anecdotal.

The whole chain from labelled pairs to the tables below is a single reproducible pipeline: splitting, baselines, training, evaluation and error analysis each read and write files, and every table in this paper is generated from those files rather than transcribed.

\begin{table*}[!t]
\centering
\caption{E1: baselines and LMP model on annotated test pairs}
\label{tab:e1}
\footnotesize
\begin{tabularx}{\textwidth}{Y>{\centering\arraybackslash}p{58pt}>{\centering\arraybackslash}p{58pt}>{\centering\arraybackslash}p{54pt}>{\centering\arraybackslash}p{54pt}>{\centering\arraybackslash}p{46pt}}
\toprule
\textbf{System} & \textbf{Pearson} & \textbf{Spearman} & \textbf{RMSE} & \textbf{MAE} & \textbf{Over-score \%} \\
\midrule
XLM-R, FrJUDGE\ensuremath{\rightarrow}ours, DA & 0.576 \ensuremath{\pm} 0.102 & 0.544 \ensuremath{\pm} 0.106 & 2.037 \ensuremath{\pm} 0.142 & 1.474 \ensuremath{\pm} 0.213 & 38.8 \ensuremath{\pm} 3.9 \\
XLM-R, FrJUDGE\ensuremath{\rightarrow}ours, no DA & 0.586 \ensuremath{\pm} 0.098 & 0.553 \ensuremath{\pm} 0.101 & 2.037 \ensuremath{\pm} 0.110 & 1.525 \ensuremath{\pm} 0.042 & 46.1 \ensuremath{\pm} 7.0 \\
XLM-R from scratch, DA & 0.228 \ensuremath{\pm} 0.288 & 0.259 \ensuremath{\pm} 0.280 & 2.464 \ensuremath{\pm} 0.224 & 2.020 \ensuremath{\pm} 0.469 & 47.9 \ensuremath{\pm} 12.1 \\
XLM-R from scratch, no DA & -0.030 \ensuremath{\pm} 0.048 & -0.020 \ensuremath{\pm} 0.044 & 2.444 \ensuremath{\pm} 0.090 & 2.245 \ensuremath{\pm} 0.079 & 45.8 \ensuremath{\pm} 5.6 \\
MuRIL from scratch, DA & 0.067 \ensuremath{\pm} 0.097 & 0.097 \ensuremath{\pm} 0.060 & 2.607 \ensuremath{\pm} 0.180 & 2.365 \ensuremath{\pm} 0.135 & 41.7 \ensuremath{\pm} 6.6 \\
MuRIL from scratch, no DA & -0.042 \ensuremath{\pm} 0.014 & -0.034 \ensuremath{\pm} 0.028 & 4.443 \ensuremath{\pm} 0.058 & 3.876 \ensuremath{\pm} 0.084 & 6.4 \ensuremath{\pm} 2.1 \\
labse-cal & 0.101 \ensuremath{\pm} 0.026 & 0.224 \ensuremath{\pm} 0.066 & 2.417 \ensuremath{\pm} 0.124 & 2.184 \ensuremath{\pm} 0.090 & 44.4 \ensuremath{\pm} 4.8 \\
labse & 0.101 \ensuremath{\pm} 0.026 & 0.224 \ensuremath{\pm} 0.066 & 2.811 \ensuremath{\pm} 0.157 & 2.110 \ensuremath{\pm} 0.121 & 41.0 \ensuremath{\pm} 3.6 \\
bertscore-cal & 0.055 \ensuremath{\pm} 0.040 & 0.115 \ensuremath{\pm} 0.033 & 2.427 \ensuremath{\pm} 0.120 & 2.206 \ensuremath{\pm} 0.088 & 44.7 \ensuremath{\pm} 4.9 \\
bertscore & 0.055 \ensuremath{\pm} 0.040 & 0.115 \ensuremath{\pm} 0.033 & 2.974 \ensuremath{\pm} 0.235 & 2.180 \ensuremath{\pm} 0.232 & 59.2 \ensuremath{\pm} 1.6 \\
bleu-cal & 0.018 \ensuremath{\pm} 0.021 & 0.088 \ensuremath{\pm} 0.042 & 2.433 \ensuremath{\pm} 0.120 & 2.212 \ensuremath{\pm} 0.087 & 45.4 \ensuremath{\pm} 5.8 \\
bleu & 0.018 \ensuremath{\pm} 0.021 & 0.088 \ensuremath{\pm} 0.042 & 5.512 \ensuremath{\pm} 0.238 & 4.977 \ensuremath{\pm} 0.316 & 6.4 \ensuremath{\pm} 1.3 \\
rougeL-cal & 0.010 \ensuremath{\pm} 0.033 & 0.071 \ensuremath{\pm} 0.040 & 2.430 \ensuremath{\pm} 0.119 & 2.217 \ensuremath{\pm} 0.090 & 45.4 \ensuremath{\pm} 5.8 \\
rougeL & 0.010 \ensuremath{\pm} 0.033 & 0.071 \ensuremath{\pm} 0.040 & 4.041 \ensuremath{\pm} 0.159 & 3.563 \ensuremath{\pm} 0.169 & 21.2 \ensuremath{\pm} 1.8 \\
Variant-mean lookup (uses answer key; diagnostic, not a system) & 0.831 \ensuremath{\pm} 0.063 & n/a & 1.385 \ensuremath{\pm} 0.138 & 0.868 \ensuremath{\pm} 0.045 & 33.7 \ensuremath{\pm} 16.7 \\
\bottomrule
\end{tabularx}
\end{table*}

\begin{table*}[!t]
\centering
\caption{E2: results by language, and zero-shot transfer}
\label{tab:e2}
\footnotesize
\begin{tabularx}{\textwidth}{YY>{\centering\arraybackslash}p{58pt}>{\centering\arraybackslash}p{54pt}>{\centering\arraybackslash}p{42pt}}
\toprule
\textbf{System} & \textbf{Setting} & \textbf{Pearson} & \textbf{RMSE} & \textbf{Over-score \%} \\
\midrule
XLM-R, FrJUDGE\ensuremath{\rightarrow}ours, DA & train EN+HI, test EN & 0.165 \ensuremath{\pm} 0.063 & 2.380 \ensuremath{\pm} 0.161 & 36.7 \ensuremath{\pm} 8.8 \\
XLM-R, FrJUDGE\ensuremath{\rightarrow}ours, DA & train EN+HI, test HI & -0.041 \ensuremath{\pm} 0.144 & 2.638 \ensuremath{\pm} 0.334 & 45.6 \ensuremath{\pm} 3.8 \\
XLM-R, FrJUDGE\ensuremath{\rightarrow}ours, DA & train EN \ensuremath{\rightarrow} test HI (zero-shot) & 0.035 \ensuremath{\pm} 0.027 & 2.564 \ensuremath{\pm} 0.293 & 42.2 \ensuremath{\pm} 1.9 \\
XLM-R, FrJUDGE\ensuremath{\rightarrow}ours, DA & train HI \ensuremath{\rightarrow} test EN (zero-shot) & 0.131 \ensuremath{\pm} 0.115 & 2.519 \ensuremath{\pm} 0.282 & 31.1 \ensuremath{\pm} 6.9 \\
XLM-R from scratch, DA & train EN+HI, test EN & -0.059 \ensuremath{\pm} 0.197 & 2.391 \ensuremath{\pm} 0.141 & 38.9 \ensuremath{\pm} 7.7 \\
XLM-R from scratch, DA & train EN+HI, test HI & -0.058 \ensuremath{\pm} 0.108 & 2.314 \ensuremath{\pm} 0.146 & 46.7 \ensuremath{\pm} 3.3 \\
XLM-R from scratch, DA & train EN \ensuremath{\rightarrow} test HI (zero-shot) & -0.059 \ensuremath{\pm} 0.086 & 2.583 \ensuremath{\pm} 0.315 & 44.4 \ensuremath{\pm} 3.8 \\
XLM-R from scratch, DA & train HI \ensuremath{\rightarrow} test EN (zero-shot) & 0.028 \ensuremath{\pm} 0.027 & 2.433 \ensuremath{\pm} 0.074 & 37.8 \ensuremath{\pm} 6.9 \\
\bottomrule
\end{tabularx}
\end{table*}

\begin{table*}[!t]
\centering
\caption{E3: cross-act generalisation}
\label{tab:e3}
\footnotesize
\begin{tabularx}{\textwidth}{YY>{\centering\arraybackslash}p{54pt}>{\centering\arraybackslash}p{54pt}>{\centering\arraybackslash}p{54pt}}
\toprule
\textbf{System} & \textbf{Split} & \textbf{Pearson} & \textbf{RMSE} & \textbf{MAE} \\
\midrule
XLM-R, FrJUDGE\ensuremath{\rightarrow}ours, DA & CPA \ensuremath{\rightarrow} RTI & 0.388 \ensuremath{\pm} 0.021 & 2.391 \ensuremath{\pm} 0.101 & 1.703 \ensuremath{\pm} 0.031 \\
XLM-R, FrJUDGE\ensuremath{\rightarrow}ours, DA & RTI \ensuremath{\rightarrow} CPA & 0.119 \ensuremath{\pm} 0.051 & 2.629 \ensuremath{\pm} 0.068 & 2.212 \ensuremath{\pm} 0.041 \\
XLM-R from scratch, DA & CPA \ensuremath{\rightarrow} RTI & 0.186 \ensuremath{\pm} 0.053 & 2.603 \ensuremath{\pm} 0.163 & 2.042 \ensuremath{\pm} 0.162 \\
XLM-R from scratch, DA & RTI \ensuremath{\rightarrow} CPA & 0.056 \ensuremath{\pm} 0.110 & 2.589 \ensuremath{\pm} 0.036 & 2.299 \ensuremath{\pm} 0.069 \\
\bottomrule
\end{tabularx}
\end{table*}

\begin{table*}[!t]
\centering
\caption{E4: sanity checks on disjoint holdouts}
\label{tab:e4}
\footnotesize
\begin{tabularx}{\textwidth}{Y>{\centering\arraybackslash}p{46pt}>{\centering\arraybackslash}p{46pt}>{\centering\arraybackslash}p{46pt}Y}
\toprule
\textbf{System} & \textbf{Identical \ensuremath{\geq} 9.9 \%: EN (n=50)} & \textbf{Identical \ensuremath{\geq} 9.9 \%: HI (n=3)} & \textbf{Unrelated \ensuremath{\leq} 1.1 \%: EN (n=17)} & \textbf{Unrelated \ensuremath{\leq} 1.1 \%: HI (n=0)} \\
\midrule
XLM-R, FrJUDGE\ensuremath{\rightarrow}ours, DA & 66.7 \ensuremath{\pm} 57.7 & 66.7 \ensuremath{\pm} 57.7 & 94.1 \ensuremath{\pm} 5.9 & not evaluated \\
XLM-R, FrJUDGE\ensuremath{\rightarrow}ours, no DA & 91.3 \ensuremath{\pm} 5.0 & 66.7 \ensuremath{\pm} 0.0 & 17.6 \ensuremath{\pm} 10.2 & not evaluated \\
XLM-R from scratch, DA & 33.3 \ensuremath{\pm} 57.7 & 33.3 \ensuremath{\pm} 57.7 & 31.4 \ensuremath{\pm} 27.8 & not evaluated \\
labse-cal & 0.0 \ensuremath{\pm} 0.0 & 0.0 \ensuremath{\pm} 0.0 & 0.0 \ensuremath{\pm} 0.0 & not evaluated \\
labse & 100.0 \ensuremath{\pm} 0.0 & 100.0 \ensuremath{\pm} 0.0 & 0.0 \ensuremath{\pm} 0.0 & not evaluated \\
bertscore & 100.0 \ensuremath{\pm} 0.0 & 100.0 \ensuremath{\pm} 0.0 & 0.0 \ensuremath{\pm} 0.0 & not evaluated \\
bleu & 100.0 \ensuremath{\pm} 0.0 & 100.0 \ensuremath{\pm} 0.0 & 11.8 \ensuremath{\pm} 0.0 & not evaluated \\
rougeL & 100.0 \ensuremath{\pm} 0.0 & 100.0 \ensuremath{\pm} 0.0 & 0.0 \ensuremath{\pm} 0.0 & not evaluated \\
\bottomrule
\end{tabularx}
\end{table*}

\begin{table}[!t]
\centering
\caption{Mean absolute error on perturbed pairs by injected error type}
\label{tab:diag}
\footnotesize
\begin{tabularx}{\columnwidth}{Y>{\centering\arraybackslash}p{21pt}>{\centering\arraybackslash}p{54pt}>{\centering\arraybackslash}p{54pt}}
\toprule
\textbf{Injected error type} & \textbf{n per test split} & \textbf{MAE XLM-R, FrJUDGE\ensuremath{\rightarrow}ours, DA} & \textbf{MAE labse-cal} \\
\midrule
exception\_drop & 4 \ensuremath{\pm} 2 & 1.673 \ensuremath{\pm} 0.474 & 2.254 \ensuremath{\pm} 0.466 \\
modality & 9 \ensuremath{\pm} 2 & 3.160 \ensuremath{\pm} 0.878 & 2.496 \ensuremath{\pm} 0.334 \\
negation & 8 \ensuremath{\pm} 2 & 1.896 \ensuremath{\pm} 0.533 & 4.629 \ensuremath{\pm} 0.531 \\
number & 5 \ensuremath{\pm} 2 & 2.274 \ensuremath{\pm} 0.466 & 2.722 \ensuremath{\pm} 0.172 \\
scope & 7 \ensuremath{\pm} 2 & 1.860 \ensuremath{\pm} 0.517 & 2.672 \ensuremath{\pm} 0.208 \\
\bottomrule
\end{tabularx}
\end{table}

\begin{table*}[!t]
\centering
\caption{Within-variant behaviour of the principal system on the E1 random splits}
\label{tab:variant}
\footnotesize
\begin{tabularx}{\textwidth}{Y>{\centering\arraybackslash}p{25pt}>{\centering\arraybackslash}p{46pt}>{\centering\arraybackslash}p{54pt}>{\centering\arraybackslash}p{58pt}>{\centering\arraybackslash}p{58pt}}
\toprule
\textbf{Variant} & \textbf{n per test split} & \textbf{Mean label} & \textbf{MAE XLM-R, FrJUDGE\ensuremath{\rightarrow}ours, DA} & \textbf{Mean signed error XLM-R, FrJUDGE\ensuremath{\rightarrow}ours, DA} & \textbf{Within-variant Pearson XLM-R, FrJUDGE\ensuremath{\rightarrow}ours, DA} \\
\midrule
normal & 33 \ensuremath{\pm} 1 & 9.96 \ensuremath{\pm} 0.02 & 1.015 \ensuremath{\pm} 0.630 & -0.995 \ensuremath{\pm} 0.640 & 0.314 \ensuremath{\pm} 0.363 \\
aggressive & 33 \ensuremath{\pm} 1 & 7.64 \ensuremath{\pm} 0.27 & 1.205 \ensuremath{\pm} 0.134 & -0.683 \ensuremath{\pm} 0.294 & 0.487 \ensuremath{\pm} 0.063 \\
perturbed & 33 \ensuremath{\pm} 1 & 5.04 \ensuremath{\pm} 0.60 & 2.203 \ensuremath{\pm} 0.057 & 1.961 \ensuremath{\pm} 0.189 & 0.584 \ensuremath{\pm} 0.047 \\
\bottomrule
\end{tabularx}
\end{table*}

We state the from-scratch result before the detail, because on its own it is a negative result and burying it would be worse than reporting it. Trained only on our corpus the JUDGEBERT-style cross-encoder does not reliably beat generic similarity. The best such configuration, XLM-R with augmentation, averages Pearson 0.228 with a standard deviation of 0.288 across seeds, the individual seeds giving 0.56, 0.08 and 0.05. Over three seeds that is roughly 1.4 standard errors from zero, so the model is not distinguishable from no correlation, nor from LaBSE's 0.101. Without augmentation XLM-R sits at -0.030; MuRIL with augmentation reaches 0.067 and without it collapses to an RMSE of 4.4. Our data do not show that the target is unlearnable, only that with about 200 training pairs per split this configuration does not learn it reliably. The FrJUDGE reproduction in Section VII and the transfer experiment below are consistent with a data-size explanation, but they do not establish it, since FrJUDGE also differs in language, domain and label distribution.

The per-language breakdown of E1 shows where the weakness lies. Trained from scratch, the model reaches 0.297 on English test pairs and 0.010 on Hindi. With French pre-training the principal system reaches 0.675 \ensuremath{\pm} 0.104 on English and 0.308 \ensuremath{\pm} 0.115 on Hindi. That comparison is confounded by act: English test pairs pool both acts while Hindi pairs are Right to Information only, and the same runs score 0.669 \ensuremath{\pm} 0.069 on Consumer Protection pairs against 0.485 \ensuremath{\pm} 0.144 on Right to Information pairs. Hindi remains the weaker language by a wide margin, and with 105 pairs of OCR-recovered text we cannot separate how much of that gap is residual OCR noise, how much is subset size, and how much is XLM-R's representation of Hindi statutory register.

E2 proper uses only the 35 Right to Information clauses present in both languages, so that English and Hindi are compared on the same provisions; each split then has about 126 training pairs from 21 clauses. On that subset neither system learns. Trained jointly on both languages, the transferred model reaches 0.165 \ensuremath{\pm} 0.063 on English and -0.041 \ensuremath{\pm} 0.144 on Hindi; zero-shot it reaches 0.035 \ensuremath{\pm} 0.027 from English to Hindi and 0.131 \ensuremath{\pm} 0.115 from Hindi to English; the from-scratch model is at or below zero throughout. E2 is therefore inconclusive about cross-lingual transfer: with a training set this small, a null result cannot distinguish a language gap from a data gap. These numbers are also lower than the per-language E1 figures above, where the Hindi pairs are trained alongside all English pairs of both acts.

Cross-act generalisation (E3) is asymmetric. Training on the Consumer Protection Act and testing on the Right to Information Act gives 0.388 \ensuremath{\pm} 0.021 with transfer and 0.186 \ensuremath{\pm} 0.053 without; the reverse direction gives 0.119 \ensuremath{\pm} 0.051 and 0.056 \ensuremath{\pm} 0.110. Augmentation partners in these splits come only from the training act, so no target-act text reaches training. The two directions differ in more than the domain: the first trains on 40 English clauses and tests on both languages, the second trains on 38 clauses in both languages and tests on English only. With a single pair of acts we cannot attribute the asymmetry to any one of these differences.

What the baselines show is worth stating alongside that. Every generic metric is close to useless on this corpus: LaBSE reaches Pearson 0.101 \ensuremath{\pm} 0.026, multilingual BERTScore 0.055 \ensuremath{\pm} 0.040, BLEU 0.018 \ensuremath{\pm} 0.021 and ROUGE-L 0.010 \ensuremath{\pm} 0.033. These are far below the 0.46 that BERTScore achieves on FrJUDGE, and we attribute the difference mainly to our corpus. A third of our pairs carry an injected legal error, so the simplification is lexically and semantically almost identical to a faithful one while meaning something different; an overlap or embedding metric is not designed to see that difference. The diagnostic table makes this explicit: LaBSE's error is largest exactly where the change is smallest on the surface, 5.896 \ensuremath{\pm} 0.372 mean absolute error on negation reversals against 3.022 \ensuremath{\pm} 0.910 on dropped exceptions (naive LaBSE mapping; the diagnostic table uses the train-calibrated mapping, so its values differ). Correlations this low could also indicate a baseline implementation fault. We did not re-run the baselines on FrJUDGE to rule that out; what we can say is that the four metrics fail in the same direction and their ranking tracks surface sensitivity as expected.

The sanity checks add a caution about calibration that we had not anticipated. In their naive mapping all four baselines score 100\% of identical holdout pairs at or above 9.9, but none places unrelated pairs at or below 1.1 (BLEU manages 11.8\%), so every generic metric fails the second check. After the train-fitted linear calibration that makes their RMSE competitive, they fail both checks at 0\%: the fit compresses predictions towards the label mean, which improves RMSE while destroying exactly the two endpoints a meaning-preservation metric must get right. Calibrated RMSE for these systems should therefore be read as a scaling artefact, not as evidence of fitness, and this is why we report both mappings. The trained models do better but none passes both. From-scratch XLM-R with augmentation places 33\% of identical holdout pairs at or above 9.9 and 31\% of unrelated pairs at or below 1.1. The transferred models split the difference in opposite directions: without augmentation, 91.3\% \ensuremath{\pm} 5.0 of identical pairs but only 17.6\% of unrelated ones; with augmentation, 94.1\% \ensuremath{\pm} 5.9 of unrelated pairs but an identical-pair rate of 66.7\% with a standard deviation of 57.7, meaning two seeds pass every identical pair and one passes none. No system we tested passes both checks, as JUDGEBERT-DA did, and the principal system's instability on the identical check is a real defect. The holdouts are almost entirely English (3 Hindi identical pairs, no Hindi unrelated pairs), so these checks say little about Hindi.

The diagnostic set is where the case for a learned metric is clearest, and the pattern is the opposite of the one we expected. Negation is by far the hardest injected error for the calibrated LaBSE baseline, whose mean absolute error is 4.629 \ensuremath{\pm} 0.531 on negation reversals against roughly 2.3 to 2.7 on everything else: a single inserted “not” moves the legal meaning from one end of the rubric to the other and moves the sentence embedding almost not at all. The principal system handles negation much better, at 1.896 \ensuremath{\pm} 0.533, and its own hardest category is modality at 3.160 \ensuremath{\pm} 0.878, where “shall” becomes “may”. This is consistent with the learned metric detecting a transformation that similarity cannot see, though because the errors are rule-generated it may partly be recognising the edit template rather than its legal effect. Its weakness on modality is the natural next target: the obligation-to-permission flip is a common real failure of legal simplification.

Overall correlation on this corpus rewards telling the three variants apart, and a lookup that predicts each test pair by the mean training label of its variant, using answer-key information no deployed system would have, reaches 0.831 \ensuremath{\pm} 0.063, above every learned model. The within-variant table is therefore the more honest measure. Within each variant the principal system still correlates positively with the labels, 0.487 \ensuremath{\pm} 0.063 on aggressive and 0.584 \ensuremath{\pm} 0.047 on perturbed pairs, so it is not merely recognising the variant; the 0.314 \ensuremath{\pm} 0.363 on normal pairs means little because their labels are almost constant (9.96). The table also shows a systematic bias: the model under-scores faithful simplifications by about one point and over-scores perturbed ones by about two, which is the over-scoring failure on exactly the pairs that matter most.

\subsection{Transfer from FrJUDGE}

Because the from-scratch model is unstable, we add an experiment the original work does not have: train the cross-encoder on FrJUDGE first, which is the same task and rubric in French insurance text and where our code reaches 0.707, then fine-tune on our Indian splits. We report the transferred systems alongside the from-scratch rows rather than replacing them, so that the contribution of the transfer is visible as an ablation: with and without augmentation, and with and without the French initialisation.

Pre-trained on FrJUDGE with augmentation, the cross-encoder reaches Pearson 0.85 to 0.87 on the FrJUDGE test set across our two pre-training runs. Fine-tuned from there with augmentation, our principal system reaches 0.564, 0.684 and 0.481 across the three seeds (0.576 \ensuremath{\pm} 0.102); without augmentation, 0.625, 0.658 and 0.474 (0.586 \ensuremath{\pm} 0.098). We designate the augmented variant as principal because it mirrors JUDGEBERT-DA and is the configuration used in E2 to E4; this designation was made after the E1 numbers were known, and the two variants are within noise of each other. Against the from-scratch model's 0.228 \ensuremath{\pm} 0.288 the gain is large in level and also in stability: the transferred seeds lie between 0.48 and 0.68, the from-scratch seeds between 0.05 and 0.56.

What this shows is that initialising from a model trained on the same rubric in another language gives a large and more stable gain on our labels. It does not show that legal meaning preservation is language- or jurisdiction-independent: French pre-training also teaches the output scale and a label distribution, the reference labels are LLM judgements, and E2 finds no cross-lingual transfer within our own corpus. Nor is the transferred model fit to deploy: the variant lookup still beats it, it fails one sanity check on one seed in three, and Section VI shows the cases it still gets badly wrong.

\subsection{Follow-up: MuRIL and five seeds}

After the experiments above we ran two follow-ups, both post hoc and reported as such. We repeated the transfer with MuRIL [6], an encoder pre-trained on Indian languages, over five clause-grouped seeds (42 to 46), and we re-ran the XLM-R transfer over the same five seeds, which required a fresh FrJUDGE pre-training run. Table IX compares them with the principal system. We leave the principal system unchanged: promoting a better one after seeing these numbers would be selection on the test data.

\begin{table*}[!t]
\centering
\caption{Follow-up: FrJUDGE\ensuremath{\rightarrow}ours with augmentation, Pearson per experiment (mean \ensuremath{\pm} std over seeds; FrJUDGE = test Pearson of the pre-training run)}
\label{tab:followup}
\footnotesize
\begin{tabularx}{\textwidth}{Y>{\centering\arraybackslash}p{54pt}>{\centering\arraybackslash}p{54pt}>{\centering\arraybackslash}p{54pt}>{\centering\arraybackslash}p{54pt}>{\centering\arraybackslash}p{58pt}>{\centering\arraybackslash}p{54pt}>{\centering\arraybackslash}p{54pt}}
\toprule
\textbf{System (FrJUDGE\ensuremath{\rightarrow}ours, DA)} & \textbf{FrJUDGE} & \textbf{E1} & \textbf{E1 HI} & \textbf{E2 EN} & \textbf{E2 HI} & \textbf{E3 C\ensuremath{\rightarrow}R} & \textbf{E3 R\ensuremath{\rightarrow}C} \\
\midrule
XLM-R, original pre-training (3 seeds) & 0.873 / 0.849 & 0.576 \ensuremath{\pm} 0.102 & 0.308 \ensuremath{\pm} 0.115 & 0.165 \ensuremath{\pm} 0.063 & -0.041 \ensuremath{\pm} 0.144 & 0.388 \ensuremath{\pm} 0.021 & 0.119 \ensuremath{\pm} 0.051 \\
XLM-R, re-run pre-training (5 seeds) & 0.691 & 0.486 \ensuremath{\pm} 0.112 & 0.290 \ensuremath{\pm} 0.068 & 0.155 \ensuremath{\pm} 0.094 & 0.066 \ensuremath{\pm} 0.141 & 0.339 \ensuremath{\pm} 0.046 & 0.147 \ensuremath{\pm} 0.040 \\
MuRIL (5 seeds) & 0.839 & 0.577 \ensuremath{\pm} 0.068 & 0.466 \ensuremath{\pm} 0.092 & 0.269 \ensuremath{\pm} 0.090 & 0.182 \ensuremath{\pm} 0.119 & 0.394 \ensuremath{\pm} 0.054 & 0.251 \ensuremath{\pm} 0.186 \\
\bottomrule
\end{tabularx}
\end{table*}

MuRIL generalises best. Its overall correlation equals the principal system's (0.577 \ensuremath{\pm} 0.068 against 0.576 \ensuremath{\pm} 0.102) with lower variance, and it is better on every generalisation measure: Hindi in E1 (0.466 against 0.308), both languages on the matched bilingual subset (0.269 and 0.182 against 0.165 and -0.041), and the harder cross-act direction (0.251 against 0.119, though with a standard deviation of 0.186). Zero-shot transfer between the languages stays near zero for every system (MuRIL 0.086 from English to Hindi, 0.104 in reverse).

The re-run exposes a weakness in the transfer recipe. The same XLM-R pre-training on FrJUDGE, with the same seed, reached 0.691 on the FrJUDGE test set instead of 0.873, and the downstream correlation followed: 0.486 \ensuremath{\pm} 0.112 with augmentation, and 0.166 \ensuremath{\pm} 0.214 without it, against the 0.586 \ensuremath{\pm} 0.098 reported for that variant in Table III. The non-augmented result therefore rests on one favourable pre-training run; the augmented variants, and MuRIL, are the robust ones.

\subsection{Legal-error type classification}

The brief asks, optionally, for the type of legal error in each pair. We frame it as six-way classification (none, modality, negation, number, exception dropped, scope) on the standard and perturbed pairs of the E1 splits. The targets are objective: “none” for standard simplifications and the injected type for perturbed ones, not the LLM annotations; aggressive pairs have no known type and are excluded. We compare a rule baseline built on the same cues as the demonstration's cue table with cross-encoder classifiers (class-weighted cross-entropy, best development macro-F1).

\begin{table*}[!t]
\centering
\caption{Legal-error type classification on the E1 splits (five seeds): macro-F1 (mean \ensuremath{\pm} std), accuracy and per-class F1}
\label{tab:errcls_main}
\footnotesize
\begin{tabularx}{\textwidth}{Y>{\centering\arraybackslash}p{54pt}>{\centering\arraybackslash}p{21pt}>{\centering\arraybackslash}p{17pt}>{\centering\arraybackslash}p{17pt}>{\centering\arraybackslash}p{17pt}>{\centering\arraybackslash}p{17pt}>{\centering\arraybackslash}p{21pt}}
\toprule
\textbf{System} & \textbf{Macro-F1} & \textbf{Acc.} & \textbf{Mod.} & \textbf{Neg.} & \textbf{Num.} & \textbf{Exc.} & \textbf{Scope} \\
\midrule
Rules (no training) & 0.509 \ensuremath{\pm} 0.074 & 0.552 & 0.45 & 0.63 & 0.58 & 0.33 & 0.46 \\
Majority (“none”) & 0.111 \ensuremath{\pm} 0.001 & 0.502 & 0.00 & 0.00 & 0.00 & 0.00 & 0.00 \\
MuRIL & 0.362 \ensuremath{\pm} 0.174 & 0.327 & 0.25 & 0.64 & 0.57 & 0.07 & 0.47 \\
MuRIL + augmentation & 0.789 \ensuremath{\pm} 0.101 & 0.754 & 0.66 & 0.85 & 0.97 & 0.79 & 0.79 \\
XLM-R + augmentation & 0.581 \ensuremath{\pm} 0.288 & 0.695 & 0.55 & 0.64 & 0.69 & 0.30 & 0.55 \\
MuRIL + augmentation + cues & 0.573 \ensuremath{\pm} 0.072 & 0.598 & 0.47 & 0.74 & 0.73 & 0.38 & 0.49 \\
\bottomrule
\end{tabularx}
\end{table*}

Trained on the corpus alone, the classifiers lose to the rules (MuRIL 0.362 \ensuremath{\pm} 0.174 against 0.509 \ensuremath{\pm} 0.074): about 130 training pairs over six classes are too few. Augmenting only the training clauses, by injecting each applicable error type into their standard simplifications and into their original texts (about 630 training pairs per split), changes that: MuRIL reaches 0.789 \ensuremath{\pm} 0.101 (0.852, 0.874, 0.621, 0.779 and 0.817 across seeds) and beats the rules on every class. Appending the rule cues to the input does not help (0.573). Because the augmentation and the rules use the same rules that injected the test errors, this shows that the model recognises these error types; it says nothing about errors a simplifier makes naturally, which carry no type label here.

\section{Error Analysis}

Scoring the corpus blind made one structure plain, and it organises this section. The ways a simplification can betray a provision fall into two families. The first is loss: a proviso, an override, a subject-to clause, the date from which a limitation period runs, or an exception simply disappears. This is what aggressive plain-language rewriting does by its nature, and the resulting text is honest about being shorter. The second is reversal: a negation is inserted or dropped, an obligation becomes a permission or the reverse, or a quantifier narrows, so that “any political party” becomes “certain political party” and “every order” becomes “some order”. Reversal is the dangerous family, because the text stays the same length and reads as fluently as a faithful version while telling the reader the opposite of the law. It is also the family that overlap and embedding metrics cannot see: as Section V reports, LaBSE's error is nearly twice as large on negation reversals as on dropped exceptions, which is precisely backwards from how much the surface changed.

The cases below are selected automatically from our labelled pairs and the blinded annotator's own workbook comments, taking the injected-error pairs first so that the error type is known rather than inferred, and capping each group so that every error type appears before any repeats. For each case the table gives the provision, the language, the variant, the kind of change, the label (the mean of the annotators), the prediction of each system with test predictions for that pair, and the blinded annotator's own score, so that annotator disagreement stays visible rather than being hidden inside the mean.

\begin{table*}[!t]
\centering
\caption{E5: legally difficult cases, with the change that was made (pert. = perturbed, aggr. = aggressive, norm. = normal; Annot. = annotator B, the blinded one)}
\label{tab:e5}
\footnotesize
\begin{tabularx}{\textwidth}{YYY>{\centering\arraybackslash}p{21pt}>{\centering\arraybackslash}p{25pt}>{\centering\arraybackslash}p{46pt}>{\centering\arraybackslash}p{33pt}>{\centering\arraybackslash}p{37pt}}
\toprule
\textbf{Provision} & \textbf{Variant} & \textbf{Change} & \textbf{Label} & \textbf{Annot.} & \textbf{XLM-R+fr+DA} & \textbf{XLM-R+DA} & \textbf{LaBSE-cal} \\
\midrule
rti-s11\_2 HI & pert. & number & 5.0 & 5 & -- & -- & -- \\
cpa-s10\_1 EN & pert. & negation & 2.5 & 1 & 3.2 & 4.8 & 7.6 \\
rti-s11\_3 HI & pert. & number & 5.0 & 5 & -- & -- & -- \\
cpa-s10\_2 EN & pert. & negation & 2.5 & 1 & 2.5 & 6.8 & 7.7 \\
rti-s11\_4 HI & pert. & modality & 5.0 & 5 & 8.5 & 8.3 & 7.7 \\
cpa-s12 EN & pert. & scope & 5.0 & 5 & 6.0 & 6.9 & 7.9 \\
rti-s12\_1 HI & pert. & modality & 5.0 & 5 & 10.0 & 8.3 & 7.9 \\
cpa-s14\_1 EN & pert. & negation & 2.5 & 1 & 2.6 & 7.1 & 7.5 \\
rti-s12\_4 HI & pert. & scope & 5.0 & 5 & 10.0 & 7.9 & 7.8 \\
cpa-s33\_1 EN & pert. & modality & 5.0 & 5 & 8.8 & 9.1 & 7.7 \\
rti-s14\_2 HI & pert. & exception drop & 7.0 & 9 & 8.0 & 9.5 & 7.7 \\
rti-s16\_1 HI & pert. & exception drop & 5.0 & 5 & -- & -- & -- \\
rti-s11\_2 HI & aggr. & omission & 6.5 & 4 & -- & -- & -- \\
rti-s11\_3 HI & aggr. & condition/ omission & 5.5 & 3 & -- & -- & -- \\
rti-s12\_6 HI & norm. & terminology & 9.5 & 9 & -- & -- & -- \\
rti-s14\_1 HI & norm. & obligation/ permission & 9.5 & 9 & 10.0 & 7.2 & 7.7 \\
rti-s15\_5 HI & aggr. & omission/ scope & 6.5 & 4 & 7.4 & 7.8 & 7.7 \\
cpa-s51\_2 EN & aggr. & exception/ omission & 8.5 & 8 & 7.9 & 9.6 & 7.2 \\
\bottomrule
\end{tabularx}
\end{table*}

{\footnotesize
\noindent\textbf{rti-s11\_2 (HI, perturbed).} Original: ``\dev{जहां उपधारा (1) के अधीन, यथास्थिति, केन्द्रीय लोक सूचना अधिकारी या राज्य लोक सूचना अधिकारी द्वारा पर व्यक्ति पर किसी सूचना या अभिलेख या उसके किसी भाग के बारे में किसी सूचना की ताम}…'' Simplification: ``\dev{जब केन्द्रीय या राज्य लोक सूचना अधिकारी उपधारा (1) के तहत किसी सूचना या अभिलेख के बारे में पर व्यक्ति को नोटिस भेजता है, तो उस पर व्यक्ति को नोटिस मिलने की तारीख से पंद्रह दिन के}…'' The third party has ten days from receipt of the notice, not fifteen.

\noindent\textbf{cpa-s10\_1 (EN, perturbed).} Original: ``The Central Government shall, by notification, establish with effect from such date as it may specify in that notification, a Central Consumer Protection Authority to be known as…'' Simplification: ``The Central Government will not set up, by notification and from a date it chooses, a Central Consumer Protection Authority, called the Central Authority. Its job is to deal with…'' Says the Central Government will NOT establish the Central Authority, reversing the provision and contradicting the rest of the sentence.

\noindent\textbf{rti-s11\_3 (HI, perturbed).} Original: ``\dev{धारा 7 में किसी बात के होते हुए भी, यथास्थिति, केन्द्रीय लोक सूचना अधिकारी या राज्य लोक सूचना अधिकारी धारा 6 के अधीन अनुरोध प्राप्त होने के पश्चात् चालीस दिन के भीतर, यदि पर व्यक्}…'' Simplification: ``\dev{धारा 7 में कुछ भी लिखा हो, अगर पर व्यक्ति को उपधारा (2) के तहत अपनी बात रखने का मौका दिया गया है, तो केन्द्रीय या राज्य लोक सूचना अधिकारी धारा 6 के तहत अनुरोध मिलने के तीस दिन के}…'' The forty-day period for the decision becomes thirty days.

\noindent\textbf{cpa-s10\_2 (EN, perturbed).} Original: ``The Central Authority shall consist of a Chief Commissioner and such number of other Commissioners as may be prescribed, to be appointed by the Central Government to exercise the…'' Simplification: ``The Central Authority will not be made up of a Chief Commissioner and a prescribed number of other Commissioners, all appointed by the Central Government to use the powers and car…'' Says the Central Authority will NOT consist of a Chief Commissioner and the prescribed Commissioners, reversing the composition provision.

\noindent\textbf{rti-s11\_4 (HI, perturbed).} Original: ``\dev{उपधारा (3) के अधीन दी गई सूचना में यह कथन भी सम्मिलित होगा कि वह पर व्यक्ति, जिसे सूचना दी गई है, धारा 19 के अधीन उक्त विनिश्चय के विरुद्ध अपील करने का हकदार है।}'' Simplification: ``\dev{उपधारा (3) के तहत भेजी गई सूचना में यह भी लिखा होगा कि जिस पर व्यक्ति को सूचना दी गई है, वह धारा 19 के तहत उस फ़ैसले के ख़िलाफ़ अपील करेगा।}'' The third party's entitlement to appeal (\dev{हकदार है}) becomes a statement that he will appeal (\dev{अपील करेगा}), turning a right into a duty.

\noindent\textbf{cpa-s12 (EN, perturbed).} Original: ``No act or proceeding of the Central Authority shall be invalid merely by reason of— (a) any vacancy in, or any defect in the constitution of, the Central Authority; or (b) any def…'' Simplification: ``No action or proceeding of the Central Authority will become invalid just because of: (a) a vacancy in, or a defect in how the Central Authority was formed; (b) a defect in the ap…'' Clause (c) covers 'any irregularity' not affecting the merits, not 'certain irregularity'.

\noindent\textbf{rti-s12\_1 (HI, perturbed).} Original: ``\dev{केन्द्रीय सरकार, राजपत्र में अधिसूचना द्वारा, केन्द्रीय सूचना आयोग के नाम से ज्ञात एक निकाय का गठन करेगी, जो ऐसी शक्तियों का प्रयोग और ऐसे कृत्यों का पालन करेगा, जो उसे इस अधिनियम}…'' Simplification: ``\dev{केन्द्रीय सरकार राजपत्र में अधिसूचना जारी करके केन्द्रीय सूचना आयोग नाम का एक निकाय बनाएगी। यह आयोग इस अधिनियम के तहत दी गई शक्तियों का इस्तेमाल करेगा और सौंपे गए काम कर सकेगा।}'' The Commission's duty to perform its assigned functions becomes optional ('\dev{सौंपे गए काम कर सकेगा}'), although the powers limb stays mandatory.

\noindent\textbf{cpa-s14\_1 (EN, perturbed).} Original: ``The Central Authority shall regulate the procedure for transaction of its business and allocation of its business amongst the Chief Commissioner and Commissioners as may be specif…'' Simplification: ``The Central Authority will not set its own procedure for carrying out its work and for dividing the work between the Chief Commissioner and the Commissioners, as laid down in regu…'' Says the Central Authority will NOT regulate its procedure and allocation of business, reversing the provision.

\noindent\textbf{rti-s12\_4 (HI, perturbed).} Original: ``\dev{केन्द्रीय सूचना आयोग के कार्यों का साधारण अधीक्षण, निदेशन और प्रबंधन, मुख्य सूचना आयुक्त में निहित होगा, जिसकी सहायता सूचना आयुक्तों द्वारा की जाएगी और वह ऐसी सभी शक्तियों का प्रय}…'' Simplification: ``\dev{केन्द्रीय सूचना आयोग के कामकाज की देखरेख, निर्देशन और प्रबंधन मुख्य सूचना आयुक्त के पास होगा, और सूचना आयुक्त उसकी मदद करेंगे। मुख्य सूचना आयुक्त आयोग की सभी शक्तियों का इस्तेमाल}…'' '\dev{ऐसे सभी कार्य}' (all such acts and things) becomes '\dev{कुछ काम}' (some acts), narrowing what the Chief Information Commissioner may do.

\noindent\textbf{cpa-s33\_1 (EN, perturbed).} Original: ``The State Government shall provide the District Commission with such officers and other employees as may be required to assist the District Commission in the discharge of its func…'' Simplification: ``The State Government may give the District Commission the officers and staff it needs to help it do its work.'' 'Shall provide' the needed officers and employees becomes 'may give', turning a duty into a discretion.

\noindent\textbf{rti-s14\_2 (HI, perturbed).} Original: ``\dev{राष्ट्रपति, उस मुख्य सूचना आयुक्त या सूचना आयुक्त को, जिसके विरुद्ध उपधारा (1) के अधीन उच्चतम न्यायालय को निर्देश किया गया है, ऐसे निर्देश पर उच्चतम न्यायालय की रिपोर्ट प्राप्त हो}…'' Simplification: ``\dev{जिस मुख्य सूचना आयुक्त या सूचना आयुक्त के ख़िलाफ़ उपधारा (1) के तहत उच्चतम न्यायालय को मामला भेजा गया है, राष्ट्रपति उसे पद से निलंबित कर सकते हैं। ज़रूरी लगे तो जाँच के दौरान उसे}…'' Drops how long the suspension lasts, namely until the President passes orders on receiving the Supreme Court's report.

\noindent\textbf{rti-s16\_1 (HI, perturbed).} Original: ``\dev{राज्य मुख्य सूचना आयुक्त .'ऐसी अवधि के लिए जो केन्द्रीय सरकार द्वारा विहित की जाए पद धारण करेगा और पुनर्नियुक्ति के लिए पात्र नहीं होगा : परन्तु कोई राज्य मुख्य सूचना आयुक्त पैंसठ}…'' Simplification: ``\dev{राज्य मुख्य सूचना आयुक्त केन्द्रीय सरकार द्वारा तय की गई अवधि तक पद पर रहेगा और उसे दोबारा नियुक्त नहीं किया जा सकेगा।}'' Drops the proviso barring the State Chief Information Commissioner from holding office after the age of sixty-five.

\noindent\textbf{rti-s11\_2 (HI, aggressive).} Original: ``\dev{जहां उपधारा (1) के अधीन, यथास्थिति, केन्द्रीय लोक सूचना अधिकारी या राज्य लोक सूचना अधिकारी द्वारा पर व्यक्ति पर किसी सूचना या अभिलेख या उसके किसी भाग के बारे में किसी सूचना की ताम}…'' Simplification: ``\dev{अगर किसी दूसरे व्यक्ति की जानकारी दी जानी है, तो उसे दस दिन में अपनी आपत्ति बताने का मौका मिलेगा।}'' Drops the information officer who serves the notice and the fact that the ten days run from receipt of that notice, so the time limit has no starting point.

\noindent\textbf{rti-s11\_3 (HI, aggressive).} Original: ``\dev{धारा 7 में किसी बात के होते हुए भी, यथास्थिति, केन्द्रीय लोक सूचना अधिकारी या राज्य लोक सूचना अधिकारी धारा 6 के अधीन अनुरोध प्राप्त होने के पश्चात् चालीस दिन के भीतर, यदि पर व्यक्}…'' Simplification: ``\dev{किसी दूसरे की जानकारी के मामले में अधिकारी को चालीस दिन में तय करना होता है कि जानकारी दी जाए या नहीं, और उस व्यक्ति को लिखकर बताना होता है।}'' Drops the override of section 7, the sub-section (2) opportunity condition and the receipt of the section 6 request as the start of the forty days.

\noindent\textbf{rti-s12\_6 (HI, normal).} Original: ``\dev{मुख्य सूचना आयुक्त या कोई सूचना आयुक्त, यथास्थिति, संसद् का सदस्य या किसी राज्य या संघ राज्यक्षेत्र के विधान-मंडल का सदस्य नहीं होगा या कोई अन्य लाभ का पद धारित नहीं करेगा या किसी}…'' Simplification: ``\dev{मुख्य सूचना आयुक्त या कोई सूचना आयुक्त संसद का या किसी राज्य या संघ राज्यक्षेत्र की विधानसभा का सदस्य नहीं होगा। वह कोई दूसरा लाभ का पद नहीं रखेगा, किसी राजनीतिक दल से नहीं जुड़ेग}…'' All disqualifications are kept, but '\dev{विधान-मंडल}' (Legislature) is narrowed to '\dev{विधानसभा}' (Legislative Assembly).

\noindent\textbf{rti-s14\_1 (HI, normal).} Original: ``\dev{उपधारा (3) के उपबंधों के अधीन रहते हुए, मुख्य सूचना आयुक्त या किसी सूचना आयुक्त को राष्ट्रपति के आदेश द्वारा साबित कदाचार या असमर्थता के आधार पर उसके पद से तभी हटाया जाएगा, जब उच्}…'' Simplification: ``\dev{उपधारा (3) के प्रावधानों के अधीन, मुख्य सूचना आयुक्त या किसी सूचना आयुक्त को साबित कदाचार या असमर्थता के आधार पर केवल राष्ट्रपति के आदेश से पद से हटाया जा सकेगा। ऐसा तभी होगा जब र}…'' '\dev{हटाया जाएगा}' becomes '\dev{हटाया जा सकेगा}', a modality shift, though the exclusive route via the President's order and the Supreme Court report is preserved.

\noindent\textbf{rti-s15\_5 (HI, aggressive).} Original: ``\dev{राज्य मुख्य सूचना आयुक्त और राज्य सूचना आयुक्त विधि, विज्ञान और प्रौद्योगिकी, समाजसेवा, प्रबंध, पत्रकारिता, जनसंपर्क माध्यम या प्रशासन और शासन में व्यापक ज्ञान और अनुभव वाले समाज}…'' Simplification: ``\dev{राज्य सूचना आयुक्त समाज के जाने-माने लोग होने चाहिए। उन्हें कानून, विज्ञान, पत्रकारिता या सरकार चलाने जैसे कामों की अच्छी समझ होनी चाहिए।}'' Drops the State Chief Information Commissioner and several qualifying fields, and turns the closed list of fields into an open 'such as' list.

\noindent\textbf{cpa-s51\_2 (EN, aggressive).} Original: ``Save as otherwise expressly provided under this Act or by any other law for the time being in force, an appeal shall lie to the National Commission from any order passed in appeal…'' Simplification: ``If a State Commission decides an appeal, you can go to the National Commission next. It will only take the case if there is an important legal question.'' Drops the saving clause for anything expressly provided under this or another law, and 'important legal question' replaces substantial question of law.

}
Two patterns in that table are worth drawing out. Where an injected negation reverses a whole operative provision the annotators split, one placing it Off-track and the other Imprecise, and the mean label of roughly 2.5 understates how misleading the text is: a reader told that the Central Government “will not” establish the authority the section establishes has been told the opposite of the law. Where the change is a quantifier or a single digit, by contrast, the annotators agreed closely, because the error is easy to see once noticed. The practical implication for a deployed metric is that the cases needing the most reliable detection are exactly the ones on which human, or in our case model, annotators are least consistent.

The predicted columns come only from the E1 random splits, averaged over the seeds in which the pair was a test item; a dash means the pair was never in an E1 test split. They overturn an expectation. The three outright negation reversals, all labelled 2.5, are where the principal system does best, at 3.2, 2.5 and 2.6, against 4.8 to 7.1 from scratch and 7.5 to 7.7 for calibrated LaBSE. Its failures lie elsewhere. Two Hindi Right to Information pairs labelled 5.0, a modality flip (rti-s12\_1) and a scope change (rti-s12\_4), are predicted at 10.0, a perfect score for a simplification that changes who must act or to whom the provision applies, and the English modality flip in cpa-s33\_1 receives 8.8. This matches the diagnostic table, where modality is the principal system's worst category. The over-scoring failure JUDGEBERT identified in generic metrics thus survives in our learned metric, concentrated on obligation and permission and on Hindi, and a good average error does not protect against it. It is also why we regard the diagnostic set, rather than any correlation we report, as the contribution most likely to outlast this paper.

\section{Reproducibility Notes}

Two defects in our own pipeline produced plausible-looking but worthless results, and both are worth reporting because neither announces itself.

The first is a tokenizer fallback. Under transformers 5.x, loading the CamemBERT-v2 checkpoint through the generic auto-tokenizer yields a byte-level tokenizer without its merge table, which splits text into single characters. A sentence pair then exceeds 650 tokens, is truncated at the maximum length, and training collapses onto predicting the label mean: the training error equals the label variance and the test correlation is slightly negative. Nothing errors out. Our loader now detects this case and falls back to the checkpoint's own tokenizer file, and we verify the token count for any new backbone.

The second is the Hindi text layer described in Section III-C, where two different extractors fail in two different ways while both returning text that looks like Hindi. The general lesson for legal NLP in Indic scripts is that an extraction step needs an explicit, quantitative acceptance test, because downstream metrics are computed happily on corrupted input.

As an end-to-end check of the training code against a published result, we converted the released FrJUDGE splits and reran the original no-augmentation setting over the same ten seeds on a single T4. The check succeeds, with a caveat worth stating precisely. Following the released code, which trains for 100 epochs and keeps the best development checkpoint, we obtain Pearson 0.707 \ensuremath{\pm} 0.042 and RMSE 1.788 \ensuremath{\pm} 0.115 against the published 0.74 \ensuremath{\pm} 0.02 and 1.72. Our Pearson is 0.033 lower, which over ten seeds each is roughly two standard errors: close enough to confirm that the implementation is the same method in the same regime, not close enough to call an exact replication. The residual gap is most plausibly explained by library versions (the tokenizer fallback described below is itself a transformers 5.x artefact), our 512-token maximum length, and seed-level nondeterminism on different hardware.

The more useful finding is that the setting matters more than the gap. Trained with early stopping at patience 5, which is what the paper's text describes, the same code gives Pearson 0.606 \ensuremath{\pm} 0.284: one seed stopped at epoch 5 and collapsed to -0.196, and the remaining nine average 0.695. The released code has that early-stopping callback commented out. A reader following the paper's description rather than its code would therefore see a wide, unstable spread and might conclude the method does not work, when the instability belongs to the stopping rule. We report both settings for that reason.

A partial augmented run (two seeds so far) is also informative about the headline number of the original work. Trained with augmentation but evaluated on the human-annotated test pairs only, it reaches Pearson 0.869 \ensuremath{\pm} 0.023 and RMSE 1.276 \ensuremath{\pm} 0.128, against the 0.97 the paper reports on a test set that includes the trivial identical and unrelated augmentation pairs. The augmented model is genuinely much stronger than the unaugmented one, and the published 0.97 is also partly an artefact of what is in the test set. This is the concrete reason our own headline metrics exclude augmentation and sanity pairs. The FrJUDGE data is CC-BY-NC-SA 4.0; we use it for this check and as the pre-training data of the transfer experiment, and we do not redistribute it.

\section{Limitations and Ethics}

\begin{itemize}
\item The LMP labels in this version are disclosed model judgements, not human legal judgements. The central claim of the predecessor work, agreement with human legal judgement, is therefore not tested here, and the planned 30-pair human validation set is the minimum needed to begin testing it.
\item The inter-LLM agreement (Pearson 0.723) is the correlation between two instances of the same model, may reflect shared pretraining rather than a well-defined construct, and is not an upper bound on predicting their mean. One of the two also wrote the simplifications and so is not blind in practice. The 68 largest disagreements were not adjudicated.
\item The benchmark encodes its construction: predicting each test pair by the mean training label of its variant, which uses answer-key information and is not a deployable system, already reaches a higher Pearson than any learned model (E1 table). Overall correlation therefore cannot show that a model assesses legal consequences rather than recognising simplification style or error templates; the within-variant table is the more honest measure.
\item No human checked the Hindi text against the source PDFs; the OCR and rule-based repairs were verified only by an automated agent spot-check. The sanity holdouts contain 3 Hindi identical pairs and no Hindi unrelated pairs, so the unrelated-pair check is English-only.
\item The corpus covers two acts and, for Hindi, one act only, because no official Hindi text of the Consumer Protection Act was available when the corpus was built (the Hindi text of the Bill as introduced has since been located and is not used here). The two Right to Information editions differ, and the Hindi text had to be recovered by OCR; both are documented and flagged, and affected provisions are dropped rather than edited.
\item The clause review that gates the corpus was performed by a rule-based agent rather than a human, and the simplifications were written by an agent rather than sampled from an API at a recorded temperature. Both are disclosed above. A human pass over the kept clauses and over a sample of simplifications is the first thing we would add.
\item The Hindi text is OCR output. Although it passes our acceptance tests and its numbering was repaired from the sequence, OCR error in legal vocabulary cannot be excluded, and a Hindi legal reader has not yet verified the kept clauses.
\item E2 is restricted to the 35 bilingual Right to Information clauses, about 126 training pairs per split, which is too few to support a conclusion about cross-lingual transfer either way.
\item The FrJUDGE pre-training is not stable: re-running it with the same seed gave 0.691 instead of 0.873 on FrJUDGE and lowered every downstream result, so single-run transfer numbers, including our non-augmented 0.586, should be read with that variance in mind. The MuRIL and error-type results are post-hoc follow-ups and are reported as such.
\item The error-type classifier is trained and tested on rule-injected errors, with augmentation from the same rules; its macro-F1 measures recognition of those five types, not of the omissions and paraphrase errors that real simplifiers make.
\item Injected errors are rule-based and may be easier to detect than errors a model makes naturally, so perturbed results are reported separately from the rest.
\item The corpus is small, and the consequences fall hardest on the development split: about 35 annotated pairs decide both when to stop training and which checkpoint to keep. We changed the learning rate, batch size and stopping rule after seeing early stopping fail on that split, so those choices are informed by the development data and the reported numbers should be read as optimistic by an unknown margin. A larger corpus with a held-out tuning split is the right fix, and the honest statement of what we have is that the protocol was adjusted once, deliberately, and is disclosed rather than tuned repeatedly against the test set.
\item Statute texts are public government documents. The corpus is a course research artefact for evaluating simplification quality; it is not a validated legal tool, and a predicted LMP score is not legal advice.
\end{itemize}

\section{Conclusion}

We built a bilingual corpus of Indian statutory simplifications designed to make the weaknesses of the predecessor work measurable: three contrasting simplification styles including a diagnostic set with an injected legal error of a known type, clause-grouped splits verified disjoint, and sanity checks on provisions used nowhere else. Trained on that corpus alone, the learned metric is unstable, at Pearson 0.228 with a standard deviation of 0.288, and not reliably better than a generic sentence-embedding baseline. Pre-trained on FrJUDGE and then fine-tuned on the same splits, it reaches 0.576 \ensuremath{\pm} 0.102 (English 0.675, Hindi 0.308), more stably across seeds. Cross-act transfer works in one direction (0.388) and barely in the other (0.119), and the 35-clause bilingual subset is too small to show cross-lingual transfer either way.

Two follow-ups sharpen this picture. MuRIL with the same FrJUDGE pre-training matches the principal system overall (0.577 \ensuremath{\pm} 0.068 over five seeds) and generalises better, to Hindi (0.466) and to the harder cross-act direction (0.251), while a re-run of the XLM-R pre-training shows that transfer numbers depend on the pre-training run. And the legal-error type can be predicted: a MuRIL classifier trained with injected-error augmentation reaches macro-F1 0.789 against 0.509 for rules, which gives the score an explanation, within the limits of rule-injected errors.

Three cautions frame these numbers. The corpus encodes its construction, so a lookup by variant already beats every learned model, and only the within-variant correlations (0.31 to 0.58) speak to whether a model reads legal content. The diagnostic set separates the learned metric from generic ones as designed, with negations handled well and obligation-to-permission flips handled badly, and no system passes both sanity checks. And the measurement apparatus is sound, since the same code reproduces the published FrJUDGE number in the released setting while showing that the setting described in the paper's prose is unstable. Two things must happen before any of this is usable. Our disclosed model annotations have to be replaced by human legal judgement, without which the central claim of this line of work remains untested whatever correlation we report. And the failure mode in Section VI has to be fixed rather than averaged away: a metric that gives a reversed obligation a perfect score is dangerous in exactly the situation it was built for, and no aggregate correlation, ours included, makes that acceptable.

\begin{thebibliography}{99}
\bibitem{ref1} D. Beauchemin, M. Albert-Rochette, R. Khoury and P.-L. Deziel, “JUDGEBERT: Assessing Legal Meaning Preservation Between Sentences,” in Proc. EMNLP, 2025, pp. 92-118.
\bibitem{ref2} T. Kew, A. Chi, L. Vasquez-Rodriguez, S. Agrawal, D. Aumiller, F. Alva-Manchego and M. Shardlow, “BLESS: Benchmarking Large Language Models on Sentence Simplification,” in Proc. EMNLP, 2023.
\bibitem{ref3} T. Zhang, V. Kishore, F. Wu, K. Q. Weinberger and Y. Artzi, “BERTScore: Evaluating Text Generation with BERT,” in Proc. ICLR, 2020.
\bibitem{ref4} F. Feng, Y. Yang, D. Cer, N. Arivazhagan and W. Wang, “Language-agnostic BERT Sentence Embedding,” in Proc. ACL, 2022.
\bibitem{ref5} A. Conneau, K. Khandelwal, N. Goyal, V. Chaudhary, G. Wenzek, F. Guzman, E. Grave, M. Ott, L. Zettlemoyer and V. Stoyanov, “Unsupervised Cross-lingual Representation Learning at Scale,” in Proc. ACL, 2020.
\bibitem{ref6} S. Khanuja, D. Bansal, S. Mehtani, S. Khosla, A. Dey, B. Gopalan, D. K. Margam, P. Aggarwal, R. T. Nagipogu, S. Dave, S. Gupta, S. C. B. Gali, V. Subramanian and P. Talukdar, “MuRIL: Multilingual Representations for Indian Languages,” arXiv preprint arXiv:2103.10730, 2021.
\bibitem{ref7} S. Doddapaneni, R. Aralikatte, G. Ramesh, S. Goyal, M. M. Khapra, A. Kunchukuttan and P. Kumar, “Towards Leaving No Indic Language Behind: Building Monolingual Corpora, Benchmark and Models for Indic Languages,” in Proc. ACL, 2023.
\bibitem{ref8} K. Krippendorff, “Computing Krippendorff's Alpha-Reliability,” Departmental Papers (ASC), Annenberg School for Communication, University of Pennsylvania, 2011.
\bibitem{ref9} T. Scialom, P.-A. Dray, S. Lamprier, B. Piwowarski, J. Staiano, A. Wang and P. Gallinari, “QuestEval: Summarization Asks for Fact-based Evaluation,” in Proc. EMNLP, 2021.
\bibitem{ref10} M. Maddela, Y. Dou, D. Heineman and W. Xu, “LENS: A Learnable Evaluation Metric for Text Simplification,” in Proc. ACL, 2023.
\bibitem{ref11} D. Beauchemin, H. Saggion and R. Khoury, “MeaningBERT: Assessing Meaning Preservation Between Sentences,” Frontiers in Artificial Intelligence, vol. 6, 2023, doi:10.3389/frai.2023.1223924.
\bibitem{ref12} S. Paul, A. Mandal, P. Goyal and S. Ghosh, “Pre-trained Language Models for the Legal Domain: A Case Study on Indian Law,” in Proc. ICAIL, 2023, pp. 187-196.
\bibitem{ref13} A. Joshi, S. Paul, A. Sharma, P. Goyal, S. Ghosh and A. Modi, “IL-TUR: Benchmark for Indian Legal Text Understanding and Reasoning,” in Proc. ACL, 2024.
\end{thebibliography}

\end{document}
